\begin{frame}{Reconfiguration problems}
    \begin{heading}{What are they?}
Ask whether one configuration can be transformed into another via a sequence of local, legal moves.
    \end{heading}
    \begin{heading}{Examples of reconfiguration problems}
    \begin{itemize}
        \item Rubik's cube
        \item Sokoban
        \item Graph recolouring reconfiguration
        \item Independent set reconfiguration
    \end{itemize}
    \end{heading}
    \pause
    Most reconfiguration problems don't belong in NP
\end{frame}

%--------------------------------------------

\begin{frame}{Independent set reconfiguration}
        Given two independent sets of size \(k\), is it possible to transform one into the other by changing vertices one by one in such a way that intermediate sets remain independent?
    \vspace{1em}
    \pause
    \begin{heading}{Two paradigms}
    \begin{itemize}
        \item Token sliding - only allows swaps within edges
        \item Token jumping - allows any swaps preserving independence
    \end{itemize}
    \end{heading}

    Both are PSPACE-complete
\end{frame}

%--------------------------------------------

\begin{frame}{Example: Token jumping reconfiguration}
\begin{figure}
\centering
\begin{tikzpicture}[scale=2]

% coordinates (defined once)
\coordinate (c1) at (1,2);
\coordinate (c2) at (3,2);
\coordinate (c3) at (0,1);
\coordinate (c4) at (2,1);
\coordinate (c5) at (4,1);
\coordinate (c6) at (1,0);
\coordinate (c7) at (3,0);

% edges (drawn once)
\draw[line width=1.5pt, dashed] (c1)--(c2);
\draw[line width=1.5pt, dashed] (c1)--(c3);
\draw[line width=1.5pt, dashed] (c1)--(c4);
\draw[line width=1.5pt, dashed] (c3)--(c6);
\draw[line width=1.5pt, dashed] (c1)--(c6);
\draw[line width=1.5pt, dashed] (c4)--(c6);
\draw[line width=1.5pt, dashed] (c2)--(c4);
\draw[line width=1.5pt, dashed] (c2)--(c5);
\draw[line width=1.5pt, dashed] (c2)--(c7);
\draw[line width=1.5pt, dashed] (c4)--(c7);
\draw[line width=1.5pt, dashed] (c5)--(c7);
\draw[line width=1.5pt, dashed] (c6)--(c7);

\only<1>{
    \node[circle,draw,fill=cfgr,minimum size=0.6cm,line width=1.5pt] at (c1) {};
    \node[circle,draw,fill=cfblue,minimum size=0.6cm,line width=1.5pt] at (c2) {};
    \node[circle,draw,fill=cfgr,minimum size=0.6cm,line width=1.5pt] at (c3) {};
    \node[circle,draw,fill=white,minimum size=0.6cm,line width=1.5pt] at (c4) {};
    \node[circle,draw,fill=cfblue,minimum size=0.6cm,line width=1.5pt] at (c5) {};
    \node[circle,draw,fill=cfgr,minimum size=0.6cm,line width=1.5pt] at (c6) {};
    \node[circle,draw,fill=cfblue,minimum size=0.6cm,line width=1.5pt] at (c7) {};
}

\only<2>{
    \node[circle,draw,fill=cfred,minimum size=0.6cm,line width=1.5pt] at (c1) {};
    \node[circle,draw,fill=white,minimum size=0.6cm,line width=1.5pt] at (c2) {};
    \node[circle,draw,fill=cfred,minimum size=0.6cm,line width=1.5pt] at (c3) {};
    \node[circle,draw,fill=white,minimum size=0.6cm,line width=1.5pt] at (c4) {};
    \node[circle,draw,fill=white,minimum size=0.6cm,line width=1.5pt] at (c5) {};
    \node[circle,draw,fill=cfred,minimum size=0.6cm,line width=1.5pt] at (c6) {};
    \node[circle,draw,fill=white,minimum size=0.6cm,line width=1.5pt] at (c7) {};
}

\only<3>{
    \node[circle, draw=black, fill=cfred, minimum size=0.6cm, line width=1.5pt] (v1) at (c1) {};
    \node[circle, draw=black, fill=white, minimum size=0.6cm, line width=1.5pt] (v2) at (c2) {};
    \node[circle, draw=black, fill=cfred!20, minimum size=0.6cm, line width=1.5pt] (v3) at (c3) {};
    \node[circle, draw=black, fill=cfred, minimum size=0.6cm, line width=1.5pt] (v4) at (c4) {};
    \node[circle, draw=black, fill=white, minimum size=0.6cm, line width=1.5pt] (v5) at (c5) {};
    \node[circle, draw=black, fill=cfred, minimum size=0.6cm, line width=1.5pt] (v6) at (c6) {};
    \node[circle, draw=black, fill=white, minimum size=0.6cm, line width=1.5pt] (v7) at (c7) {};
}

\only<4>{
    \node[circle, draw=black, fill=cfred!20, minimum size=0.6cm, line width=1.5pt] (v1) at (c1) {};
    \node[circle, draw=black, fill=white, minimum size=0.6cm, line width=1.5pt] (v2) at (c2) {};
    \node[circle, draw=black, fill=white, minimum size=0.6cm, line width=1.5pt] (v3) at (c3) {};
    \node[circle, draw=black, fill=cfred, minimum size=0.6cm, line width=1.5pt] (v4) at (c4) {};
    \node[circle, draw=black, fill=white, minimum size=0.6cm, line width=1.5pt] (v5) at (c5) {};
    \node[circle, draw=black, fill=cfred, minimum size=0.6cm, line width=1.5pt] (v6) at (c6) {};
    \node[circle, draw=black, fill=cfred, minimum size=0.6cm, line width=1.5pt] (v7) at (c7) {};
}

\only<5>{
    \node[circle, draw=black, fill=white, minimum size=0.6cm, line width=1.5pt] (v1) at (c1) {};
    \node[circle, draw=black, fill=cfred, minimum size=0.6cm, line width=1.5pt] (v2) at (c2) {};
    \node[circle, draw=black, fill=white, minimum size=0.6cm, line width=1.5pt] (v3) at (c3) {};
    \node[circle, draw=black, fill=cfred, minimum size=0.6cm, line width=1.5pt] (v4) at (c4) {};
    \node[circle, draw=black, fill=white, minimum size=0.6cm, line width=1.5pt] (v5) at (c5) {};
    \node[circle, draw=black, fill=cfred!20, minimum size=0.6cm, line width=1.5pt] (v6) at (c6) {};
    \node[circle, draw=black, fill=cfred, minimum size=0.6cm, line width=1.5pt] (v7) at (c7) {};
}

\only<6>{
        \node[circle, draw=black, fill=white, minimum size=0.6cm, line width=1.5pt] (v1) at (c1) {};
        \node[circle, draw=black, fill=cfred, minimum size=0.6cm, line width=1.5pt] (v2) at (c2) {};
        \node[circle, draw=black, fill=white, minimum size=0.6cm, line width=1.5pt] (v3) at (c3) {};
        \node[circle, draw=black, fill=cfred!20, minimum size=0.6cm, line width=1.5pt] (v4) at (c4) {};
        \node[circle, draw=black, fill=cfred, minimum size=0.6cm, line width=1.5pt] (v5) at (c5) {};
        \node[circle, draw=black, fill=white, minimum size=0.6cm, line width=1.5pt] (v6) at (c6) {};
        \node[circle, draw=black, fill=cfred, minimum size=0.6cm, line width=1.5pt] (v7) at (c7) {};
}

\end{tikzpicture}
\end{figure}
\textbf{Note}: All the figures in the talk have dashed lines as non-edges in the graph, edges are not drawn for clarity.
\end{frame}

%--------------------------------------------

\begin{frame}{}
    \begin{fact}{Reconfiguration graph \(R_k(G)\)} is a graph whose vertices correspond to independent sets of size \(k\) in \(G\). Edges correspond to one-step reconfigurations. Possibly exponential in size!
    \end{fact}
    \vspace{1em}
    \onslide<2>{
        \begin{fact}{Reconfiguration path} is a path in \(R_k(G)\).
        \end{fact}
    }
    \vspace{-2em}
    \begin{figure}
        \centering
        \begin{tikzpicture}[scale=2]
    \coordinate (v1) at (1,2);
    \coordinate (v2) at (3,2);
    \coordinate (v3) at (0,1);
    \coordinate (v4) at (2,1);
    \coordinate (v5) at (4,1);
    \coordinate (v6) at (1,0);
    \coordinate (v7) at (3,0);

    \coordinate (r1) at (1.3, 1);
    \coordinate (r2) at (0.5, 1);
    \coordinate (r3) at (2, 1.6);
    \coordinate (r4) at (2.6, 1);
    \coordinate (r5) at (3.5, 1);
    \coordinate(r6) at (2, 0.3);
    
    \draw[line width=1.5pt, cfblue] (r1) -- (r2);
    \draw[line width=1.5pt, cfblue] (r1) -- (r3);
    \draw[line width=1.5pt, cfblue] (r1) -- (r6);
    \draw[line width=1.5pt, cfblue] (r3) -- (r4);
    \draw[line width=1.5pt, cfblue] (r4) -- (r5);
    \draw[line width=1.5pt, cfblue] (r4) -- (r6);
    
    \draw[line width=1.5pt, dashed] (v1) -- (v2);
    \draw[line width=1.5pt, dashed] (v1) -- (v3);
    \draw[line width=1.5pt, dashed] (v1) -- (v4);
    \draw[line width=1.5pt, dashed] (v3) -- (v6);
    \draw[line width=1.5pt, dashed] (v1) -- (v6);
    \draw[line width=1.5pt, dashed] (v4) -- (v6);
    \draw[line width=1.5pt, dashed] (v2) -- (v4);
    \draw[line width=1.5pt, dashed] (v2) -- (v5);
    \draw[line width=1.5pt, dashed] (v2) -- (v7);
    \draw[line width=1.5pt, dashed] (v4) -- (v7);
    \draw[line width=1.5pt, dashed] (v5) -- (v7);
    \draw[line width=1.5pt, dashed] (v6) -- (v7);


    \node[circle,draw,fill=white,minimum size=0.6cm,line width=1.5pt] at (v1) {};
    \node[circle,draw,fill=white,minimum size=0.6cm,line width=1.5pt] at (v2) {};
    \node[circle,draw,fill=white,minimum size=0.6cm,line width=1.5pt] at (v3) {};
    \node[circle,draw,fill=white,minimum size=0.6cm,line width=1.5pt] at (v4) {};
    \node[circle,draw,fill=white,minimum size=0.6cm,line width=1.5pt] at (v5) {};
    \node[circle,draw,fill=white,minimum size=0.6cm,line width=1.5pt] at (v6) {};
    \node[circle,draw,fill=white,minimum size=0.6cm,line width=1.5pt] at (v7) {};

    \only<1>{
        \node[circle, draw=cfblue, fill=cfblue, minimum size=0.3cm, line width=1.5pt] at (r1) {};
        \node[circle, draw=cfblue, fill=cfblue, minimum size=0.3cm, line width=1.5pt] at (r2) {};
        \node[circle, draw=cfblue, fill=cfblue, minimum size=0.3cm, line width=1.5pt] at (r3) {};
        \node[circle, draw=cfblue, fill=cfblue, minimum size=0.3cm, line width=1.5pt] at (r4) {};
        \node[circle, draw=cfblue, fill=cfblue, minimum size=0.3cm, line width=1.5pt] at (r5) {};
        \node[circle, draw=cfblue, fill=cfblue, minimum size=0.3cm, line width=1.5pt] at (r6) {};
    }

    \only<2>{
        \node[circle, draw=cfblue, fill=cfred, minimum size=0.3cm, line width=1.5pt] at (r1) {};
        \node[circle, draw=cfblue, fill=cfred, minimum size=0.3cm, line width=1.5pt] at (r2) {};
        \node[circle, draw=cfblue, fill=cfblue, minimum size=0.3cm, line width=1.5pt] at (r3) {};
        \node[circle, draw=cfblue, fill=cfred, minimum size=0.3cm, line width=1.5pt] at (r4) {};
        \node[circle, draw=cfblue, fill=cfred, minimum size=0.3cm, line width=1.5pt] at (r5) {};
        \node[circle, draw=cfblue, fill=cfred, minimum size=0.3cm, line width=1.5pt] at (r6) {};
    }
\end{tikzpicture}

    \end{figure}
\end{frame}

%--------------------------------------------

\begin{frame}
\vspace{-3em}
    \begin{figure}
        \centering
        \begin{tikzpicture}[scale=2]
    \coordinate (v1) at (1,2);
    \coordinate (v2) at (3,2);
    \coordinate (v3) at (0,1);
    \coordinate (v4) at (2,1);
    \coordinate (v5) at (4,1);
    \coordinate (v6) at (1,0);
    \coordinate (v7) at (3,0);

    \coordinate (r1) at (1.3, 1);
    \coordinate (r2) at (0.5, 1);
    \coordinate (r3) at (2, 1.6);
    \coordinate (r4) at (2.6, 1);
    \coordinate (r5) at (3.5, 1);
    \coordinate(r6) at (2, 0.3);
    
    \draw[line width=1.5pt, dashed] (v1) -- (v2);
    \draw[line width=1.5pt, dashed] (v1) -- (v3);
    \draw[line width=1.5pt, dashed] (v1) -- (v4);
    \draw[line width=1.5pt, dashed] (v3) -- (v6);
    \draw[line width=1.5pt, dashed] (v1) -- (v6);
    \draw[line width=1.5pt, dashed] (v4) -- (v6);
    \draw[line width=1.5pt, dashed] (v2) -- (v4);
    \draw[line width=1.5pt, dashed] (v2) -- (v5);
    \draw[line width=1.5pt, dashed] (v2) -- (v7);
    \draw[line width=1.5pt, dashed] (v4) -- (v7);
    \draw[line width=1.5pt, dashed] (v5) -- (v7);
    \draw[line width=1.5pt, dashed] (v6) -- (v7);

    \draw[line width=1.5pt, cfblue] (r1) -- (r2);
    \draw[line width=1.5pt, cfblue] (r1) -- (r3);
    \draw[line width=1.5pt, cfblue] (r1) -- (r6);
    \draw[line width=1.5pt, cfblue] (r3) -- (r4);
    \draw[line width=1.5pt, cfblue] (r4) -- (r5);
    \draw[line width=1.5pt, cfblue] (r4) -- (r6);

    \node[circle,draw,fill=white,minimum size=0.6cm,line width=1.5pt] at (v1) {};
    \node[circle,draw,fill=white,minimum size=0.6cm,line width=1.5pt] at (v2) {};
    \node[circle,draw,fill=white,minimum size=0.6cm,line width=1.5pt] at (v3) {};
    \node[circle,draw,fill=white,minimum size=0.6cm,line width=1.5pt] at (v4) {};
    \node[circle,draw,fill=white,minimum size=0.6cm,line width=1.5pt] at (v5) {};
    \node[circle,draw,fill=white,minimum size=0.6cm,line width=1.5pt] at (v6) {};
    \node[circle,draw,fill=white,minimum size=0.6cm,line width=1.5pt] at (v7) {};

    \node[circle, draw=cfblue, fill=cfred, minimum size=0.3cm, line width=1.5pt] at (r1) {};
    \node[circle, draw=cfblue, fill=cfred, minimum size=0.3cm, line width=1.5pt] at (r2) {};
    \node[circle, draw=cfblue, fill=cfblue, minimum size=0.3cm, line width=1.5pt] at (r3) {};
    \node[circle, draw=cfblue, fill=cfred, minimum size=0.3cm, line width=1.5pt] at (r4) {};
    \node[circle, draw=cfblue, fill=cfred, minimum size=0.3cm, line width=1.5pt] at (r5) {};
    \node[circle, draw=cfblue, fill=cfred, minimum size=0.3cm, line width=1.5pt] at (r6) {};
\end{tikzpicture}

    \end{figure}
    \vspace{-4em}
    \begin{fact}{\(D(n,k)\)}
        is the maximum diameter of \(R_k(G)\) among all graphs on \(n\) vertices.
    \end{fact}
    \vspace{1em}
    % \begin{fact}{Diameter path} is the path realising \(D(n,k)\).
    % \end{fact}
    \pause
    \begin{center}
        \large \textcolor{cfred}{How can we bound \(D(n,k)\)?} 
    \end{center}
    We want to find a graph such that its reconfiguration diameter is the longest possible.
\end{frame}

%--------------------------------------------

\begin{frame}{Why is this important?}
    \begin{itemize}
        \item Set reconfiguration has been heavily studied
        \item No one before has studied set reconfiguration in an extremal way 
        \item Extremal reconfiguration problems have only been studied for recolouring \framemark.
    \end{itemize}
    \vspace{1em}
    \pause
    Most of the techniques used are similar to the ones used to bound the \((6,3)\) problem.
    
    \begin{heading}{(6,3) Problem}
    (Ruzsa–Szemerédi problem) asks for asks for the maximum number of hyperedges in a (6, 3)-free n-vertex 3-uniform hypergraph \framemark..
    \end{heading}
\framereference{Bonsma, Cereceda, \emph{Finding Paths between graph colourings: PSPACE-completeness and superpolynomial distances, 2009
}}
\vspace{-1em}
\framereference{I. Z. Ruzsa and E. Szemerédi. Triple systems with no six points carrying three triangles.
Combinatorics (Keszthely, 1976), Coll. Math. Soc. J. Bolyai, 18:939–945, 1978.}
\end{frame}